% The next line makes it so it compiles the main file if I hit the compile button
% when editing this file in TeXworks.
% !TEX root = ../AIDMA.tex
\chapter{Proof Methods}\label{ch:proofs}

The ability to write proofs is important. 
Although you may not find yourself writing proofs on a regular
basis in your future career, you will certainly need to make
arguments based on evidence and logic. Proof writing is
exactly that, although it is typically more formal, and the 
subjects of our proofs are generally mathematical. Nevertheless, 
what you learn here is definitely applicable way beyond writing
mathematical proofs of simple mathematical results
(which will be the focus of our proof writing).

In the context of discrete mathematics and algorithms, 
proofs will be important in several places. I will highlight just
one here: algorithms. 
When you write an algorithm it is important that the algorithm performs as expected, both in terms of
producing the correct answer and doing so quickly.  
That is, proofs are necessary in {\em algorithm correctness} and {\em algorithm analysis}. Although we will see examples of
both, we will spend more time learning about algorithm analysis.

In this chapter we will introduce you to the basics of mathematical proofs.  Along the way we will 
review some  mathematical concepts/definitions you have probably already seen, and introduce
you to some new ones that we will find useful as we proceed.  We will continue to write proofs 
and learn more advanced proof techniques as the book continues.

\section{Direct Proofs}\label{section:dProofs}
A {\em direct proof}\index{proof!direct}\index{direct proof}
is one that follows from the definitions.  
Facts previously learned help many a time when writing a direct proof.
We will begin by seeing some direct proofs about something you should already be very familiar with:
even and odd integers.

\begin{df}
Recall that:
\begin{itemize}\index{even}\index{odd}\index{parity}
\item An {\bf even integer}\index{even} is one of the form $2k$, where $k$ is an integer.
\item An {\bf odd integer}\index{odd} is one of the form $2k + 1$ where $k$ is an integer.
\item Two integers have the same {\bf parity}\index{parity} if they are both even or both odd.
\end{itemize}
\end{df}

\begin{exa}
Use the definition of even to prove that the sum of two even integers is even.
\begin{pf}
If $x$ and $y$ are even, then $x=2a$ and $y=2b$ for some integers $a$ and $b$. 
Then $x+y=2a + 2b = 2(a + b)$, which is even since $a+b$ is an integer.
\end{pf}
\end{exa}

\begin{exa}
Use the definitions of even and odd to prove that the sum of an even integer and an
odd integer is odd.
\begin{pf}
Let $a$ be an even integer and $b$ be an odd integer.
Then $a=2f$ and $b=2g + 1$ for some integers $f$ and $g$.
Then $a+b=2f + (2g + 1) = 2(f+g) + 1$. Since $f + g$
is an integer, $a+b$ is an odd integer.
\end{pf}
\end{exa}

%\begin{rem}
%The next example is the first of many {\bf Fill in the details} exercises in 
%which {\bf you} need to supply some of the details.   
%After you have filled in the blanks, compare your answers with the solutions.  
%The answers are given with semicolons (;) separating the blanks.  
%\end{rem}

\begin{fib}
Use the definitions of even and odd to prove that the sum of two odd integers is even.
\begin{pf}
If $x$ and $y$ are odd, then  $x=2c + 1$ and $y=\blankline{.75}$ for some integers $c$ and $d$. 
Then $x+y = 2c + 1 +2d + 1 = 2(c+ d+1)$.  Now $\blankline{.75}$ is an integer, so $2(c + d+1)$
is an \blankline{.75} integer. 
\end{pf}
\begin{answer}
$2d+1$;
$c+d+1$;
even
\end{answer}
\end{fib}

%\begin{rem}
%Did you notice the $\star$ in the heading of the previous example?  
%This indicates that a solution is provided.
%If you are reading the PDF file, clicking on the $\star$ will take you to the solution.
%Clicking on the number in the solution will take you back. 
%
%If you are reading the PDF, go to the back of the book to find the solutions.
%\end{rem}

\begin{exa}
Use the definitions of even and odd to prove that the product of two odd integers is odd.
\begin{pf}
Let $a$ and $b$ be odd integers.  Then
$a=2l+1$ and $b=2m +1$ for some integers $l$ and $m$. 
Then $a\cdot b= (2l + 1)(2m + 1) = 4ml + 2l + 2m + 1 =
2(2ml + l + m) + 1$
which is odd since $2ml + m +l$ is an integer.
\end{pf}
\end{exa}

\begin{fib}
Use the definitions of even and odd to prove that the product of an even integer and an odd integer is even.
\begin{pf}
Let $a$ be an even integer and $b$ be an odd integer.
Then $a=\blankline{.75}$  and $b=\blankline{.75}$ for \blankline{1.5}.
Given that, we can see that $a\cdot b= (2n)(2o+1) = \blankline{2}.$
Since \blankline{.75} is an integer, $a\cdot b$ is \blankline{1}.
\end{pf}
\begin{answer}
$2n$;
$2o+1$;
some integers $n$ and $o$;
$4no+2n=2(2no+n)$ or $2(n(2o+1))$.  (Your steps might vary slightly, but you should end up with
either $2(2no+n)$ or  $2(n(2o+1))$ in the final step);
$2no + 1$  or $n(2o+1)$;
`an even integer' or `even'.

\end{answer}
\end{fib}

These examples may seem somewhat ridiculous since they are proving such obvious facts.  However, keep in
mind that our goal is to learn techniques for writing proofs.  As we proceed the proofs will become
more complicated, but we will continue to follow the same basic techniques we are using here.  
In other words, the fact that we are proving facts about even and odd integers is not at all important.
What is important are the techniques we are learning in the process.

You may be asking yourself ``why are we wasting our time proving such obvious results''?  
If so, ask yourself this:  Would you rather be asked to prove more complicated things right away?

Think about how you learned to read and write.
You started by reading books that only had a few simple words.
As you progressed, the books and the words in them got longer.  The vocabulary increased.
You encountered increasingly complex sentence and paragraph structures.  The same is true
when you learned to write.  You began by writing the letters of the alphabet.  Then you learned
to write words, followed by sentences, paragraphs, and eventually essays.

Learning to read and write proofs follows the same procedure.  In order to know how to write 
correct proofs you first need to see some examples of them.  
But you need to go beyond just seeing them--you need to {\em understand} them.  That is the goal of
examples like the previous one.  Your brain needs to be engaged with the material as you work through
the book.  You {\em must} work through all of the examples in order to get the most out of this book.

%\begin{rem}
%Next you will see the first of many {\bf Exercises}.
%These give you an opportunity to solve a problem from start to finish 
%and then check your answer with the solution provided.
%It is important that you try each of these on your own before looking at the solution.
%You will not get as much out of the book if you skip these or jump straight to the answer
%without trying them yourself.
%\end{rem}

\begin{exer}
Use the definition of even to prove that the product of two even integers is even.\\
Proof:\\
\vspace*{2.4in}
\begin{answer}
Let $a$ and $b$ be even integers.  Then $a=2m$ and $b=2n$ for some integers $m$ and $n$.
Their product is $ab=(2m)(2n)=2(2mn)$ which is even since $2mn$ is an integer. {\em Notice that we used two different letters here!
You cannot assume $a=2n$ and $b=2n$  because then you are
assuming that $a=b$ whether or not you realize it!}
\end{answer}
\end{exer}
%
%\begin{rem}
%The next example is an {\bf Evaluate} example.  
%These examples give a problem and then provide one or more solutions to the problem 
%based on previous student solutions.
%Your job is to evaluate each solution by finding any mistakes.  
%Mistakes include not only incorrect algebra and logic, but also unclear presentation, 
%skipped steps, incorrect assumptions, over-simplification, etc.
%When you come across these examples you should write down every error you can find.  Once you are pretty sure you
%know all of the problems (if there are any), compare your evaluation to the one given in the solutions.
%\end{rem}

\begin{eval}
Evaluate the following proof that supposedly
uses the definition of odd to prove that the product of two odd integers is odd.
\begin{spf}
By definition of odd numbers, 
let $a$ be an odd integer 
$2n+1$ 
let $b$ be an odd integer 
$2q+1$.
Then $(2n+1)(2q+1)=4nq+2n+1=2(2nq+1)+1$.
Since $2nq+1$ is an integer, $2(2nq+1)+1$ is an odd integer by definition of odd.
\end{spf}
\evallinetwo

\begin{answer}
Here are my comments on the proof.
\begin{itemize}
\item
The first sentence is phrased weird--we are not letting $a$ and $b$ be odd {\em by} the definition of odd.
We are {\em using} the definition.  
\item
It does not state that $n$ and $q$ need to be integers.
\item
Although it is not incorrect, using $n$ and $q$ is just weird in this context.  
It is customary to use adjacent letters, like $n$ and $m$, or $q$ and $r$.
\item
Given the above problems, I would rephrase the first sentence as 
{\em `Let $a$ and $b$ be an odd numbers.  
Then $a=2n+1$ and $b=2m+1$ for some integers $n$ and $m$.'}
\item
There is an algebra mistake.  The product should be $2(2nq+q+n)$.
\item
If you replace $2nq+1$ with $2nq+q+n$ (twice) in the last sentence (see the previous item)
it would be a perfect finish to the proof.
\end{itemize}
\end{answer}
\end{eval}

Sometimes students get frustrated because they think that too many details are required in a proof.  
Why are mathematicians such sticklers on the details?  The next example is the first of many
that will try to demonstrate why the seemingly little details matter.

%\begin{rem}
%The {\bf Question} examples are similar to the {\bf Evaluate} ones except that they ask a specific question.
%Write down your answer in the space provided and then compare your answer with the one in the solutions.
%\end{rem}

\begin{ques}
What is wrong with the following ``proof'' that the sum of an even and an odd number is even?
\begin{spf}
Let a=2n be an even integer and $b=2m+1$ be an odd integer.  
Then $a+b=2n+2m+1=2(n+m+1/2)$.  Since we wrote $a+b$ as a multiple of $2$, it is even.
Therefore the sum of an even and an odd number is even.
\end{spf}
\answerlinetwo
\begin{answer}
Hopefully it is clear to you that the proof {\em can't} be correct since the sum of an even and an odd number is odd, not even.
The algebra is correct.  
{\em The problem is that $n+m+1/2$ is not an integer.}
In order to be even, a number must be expressed 
in the form $2k$ {\bf\em where $k$ is an integer}.  {\em Any} number can be written as $2x$ if we don't require that $x$ be an integer,
so you {\em cannot} say that a number is even because it is of the form $2x$ unless $x$ is an integer.
\end{answer}
\end{ques}

We will find the following definitions useful throughout the book.

\begin{df}\index{divides}\index{factor}\index{multiple}
\index{divisor}
\index{a fdiv@$\vert$ (divides)}
Let $b$ and $a$ be integers with $a\neq 0$.  We say that 
$b$ is {\bf divisible by} $a$ if there exists an integer $c$ such
that $b = ac$. 
If $b$ is divisible by $a$, 
we also say that $b$ is a {\bf multiple} of $a$,
$a$ is a {\bf factor} or {\bf divisor} of $b$, and that 
$a$ {\bf divides} $b$, written as $a|b$.  If $a$ does not divide $b$, we write $a\nmid b$. 
\end{df}

\vspace*{-.1in}

\begin{exa}
Since $6=2\cdot 3$, $2|6$, and $3|6$.  But $4\nmid 6$ since we cannot write $6=4\cdot c$ for any integer $c$.
\end{exa}

\vspace*{-.1in}

\begin{exa}
Prove that the product of two even integers is
divisible by $4$.
\begin{pf}
Let  $2h$ and
$2k$ be even integers. Then $(2h)(2k) = 4(hk)$. Since $hk$ is an
integer, $4(hk)$ is  divisible by $4$. 
\end{pf}
\end{exa}

\vspace*{-.1in}

\begin{fib} Prove that if $x$ is an integer and
$7$ divides $3x + 2$, then $7$ also divides $15x^2 - 11x - 14$.\\[-.35in]
\begin{pf}
Since $7$ divides $3x+2$, we know that  $3x + 2 = 7a$, where $a$ is \blankline{1.25}.
Notice that 
\begin{eqnarray*}
15x^2 - 11x - 14&=& (\blankline{.75})(\blankline{.75})\\
&=&\blankline{.5}a(5x-7).
\end{eqnarray*}
Therefore \blankline{2}.
\end{pf}
\begin{answer}
$a$ an integer; $(3x + 2)$; $(5x - 7)$;$7$; $7$ divides $15x^2 - 11x - 14$.
\end{answer}
\end{fib}

\vspace*{-.1in}

\begin{exa}
Let $a$ and $b$ be integers such that $a|b$ and $b|a$.
Prove that either $a=b$ or $a=-b$.
\begin{pf}
If $a|b$, we can write $b=ac$ for some integer $c$.
Similarly, if $b|a$, we can write $a=bd$ for some integer $d$.
Then we can write $b=ac=(bd)c$.  Dividing both sides by $b$
(which is legal, since $b|a$ implies $b\neq 0$), we
can see that $cd=1$. Since $c$ and $d$ are integers, we know
that either $c=d=1$ or $c=d=-1$.  In the first case, we
have that $a=b$, and in the second case, we have that $a=-b$.
\end{pf}
\end{exa}

\vspace*{-.1in}

\begin{eval}
Prove that if $n$ is an integer, then $n^3 - n$ is divisible by $6$.
\begin{spf}
We have $n^3 - n = (n-1)n(n + 1)$, the product of three consecutive integers.
Among three consecutive integers at least one is even
and exactly one is divisible by $3$. Since $2$ and
$3$ do not have common factors, $6$ divides the quantity
$(n-1)n(n+1)$, and so $n^3 - n$ is divisible by $6$.
\end{spf}
\evalline
\begin{answer}
This proof is correct.  Not all of the Evaluate problems have an error!
\end{answer}
\end{eval}


\begin{df}
\index{prime}
\index{composite}
A positive integer $p>1$ is {\bf prime} if its only positive factors are $1$ and $p$. 
A positive integer $c>1$ which is not prime is said to be
{\bf composite}.
\end{df}


\begin{eval}
Prove or disprove that if $a$ is a positive even integer, then it is composite.
\begin{spf}
Let $a$ be an even number.  By definition of even, $a=2k$ for some integer $k$.
Since $a>0$, clearly $k>0$.  Since $a$ has at least two factors, $2$ and $k$, $a$ is composite.
\end{spf}
\evallinetwo
\begin{answer}
The number $2$ is positive and even but is clearly not composite since it is prime.  
Since the statement is false the proof must be incorrect.  
So where is the error?  It is in the final statement.  Although $a$ can be written
as the product of $2$ and $k$, what if $k=1$ (that is, $a=2$).  In that case 
we have not demonstrated that $a$ has a factor other than $a$ or $1$, so we can't be sure
that it is composite.
\end{answer}
\end{eval}

\begin{rem}
Notice that according to the definitions given above, $1$ is neither prime nor composite.
This is one of the many things that makes $1$ special.
\end{rem}

\begin{exer}
Prove that $2$ is the only even prime number.\\
\pflinethree
\begin{answer}
If you didn't get, try this hint before reading the rest of the solution:
 Assume $a$ is an even number other than $2$ and prove that $a$ is composite.\\
Let $a>2$ be an even integer.  Then $a=2k$ for some integer $k$.  Since $a\neq 2$, 
$a$ has a factor other than $a$ or $1$.  Therefore $a$ is not prime.
Therefore $2$ is the only even prime number.
\end{answer}
\end{exer}

\begin{ques}
Did you notice that the proof in the solution to the previous exercise (you read it, right?)
did not consider the case of $0$ or negative even numbers.
Was that O.K.? Explain why or why not.\\
\answerline
\begin{answer}
It was O.K. because according to the definition of prime, only positive integers can be prime.
Therefore we only needed to consider positive even integers.
\end{answer}
\end{ques}

\begin{df}
\index{factorial}
\index{a fact@"! (factorial)}
For a non-negative integer $n$, the quantity $n!$ (read {\bf ``$n$ factorial''}) is defined as follows.
$0! = 1$ and if $n>0$ then $n!$ is the product of all the integers from $1$ to $n$ inclusive: $$n! = 1\cdot 2 \cdots n.  $$
\end{df}

\begin{exa}
$3! = 1\cdot 2 \cdot 3 = 6$, and 
$5! = 1\cdot 2 \cdot 3 \cdot 4 \cdot 5 = 120$. 
\end{exa}

\begin{exa}
Prove that if $n>0$, then $n!\leq n^n$.
\begin{pf}
Since $1\leq n$, $2\leq n$, $\cdots$, and $(n-1)\leq n$, it is easy to see that
\begin{eqnarray*}
n!&=&1\cdot 2\cdot 3\cdots n\\
&\leq&n\cdot n\cdot n\cdots n\\
&=&n^n.
\end{eqnarray*}
\end{pf}
\end{exa}

\begin{eval}
Prove that if $n>4$ is composite, then $n$ divides $(n-1)!$.
\begin{spf}
Since $n$ is composite, $n=ab$ for some integers $1<a<n-1$ and $1<b<n-1$. 
By definition of factorial, $a | (n-1)!$ and $b |(n-1)!$.
Therefore $n=ab$ divides $(n-1)!$
\end{spf}
\evallinetwo
\begin{answer}
This one has a combination of two subtle errors.
First of all, if $a|c$ and $b|c$, that does not necessarily imply that $ab|c$.
For instance, $6|12$ and $4|12$, but it should be clear that $6\cdot 4\nmid 12$.
Second, what if $a=b$?
We'll see how to fix the proof in the next example.
\end{answer}
\end{eval}

Since the previous proof wasn't correct, let's fix it.

\begin{exa}\label{ndivfactExer}
Prove that if $n>4$ is composite, then $n$ divides $(n-1)!$.
\begin{pf} 
If $n$ is not a perfect square, then we can write $n=ab$ for some integers $a$ and $b$ with $1<a<b<n-1$.
Thus, $(n-1)! = 1\cdots a\cdots b\cdots (n-1)$.  Since $a$ and $b$ are distinct numbers on the factor list,
$n=ab$ is clearly a factor of $(n-1)!$.  

If $n$ is a perfect square, then $n=a^2$ for some integer $2<a<n-1$.
Since $a>2$, $2a<a^2=n$.  Thus, $2a<n$, so
$(n-1)! = 1\cdots a\cdots 2a\cdots (n-1)$.  Then $a(2a)=2n$ is a factor of $(n-1)!$, which means
that $n$ is as well.
\end{pf}

\end{exa}

\begin{ques}
Why was it O.K. to assume $1<a<b<n-1$ in the previous proof?\\
\answerlinethree
\begin{answer}
Since $n$ is not a perfect square, we know that $a\neq b$.  
Therefore  $a<b$ or $b<a$.  Since $a$ and $b$ are just labels for
two factors of $n$, it doesn't matter which one is larger.  So we can
just assume $a$ is the smaller one without any loss of generality.
By definition of composite, we know that $a>1$.
Finally, it should be pretty clear that $b<n-1$ since if $b=n-1$, then
$n=ab=a(n-1)\geq 2(n-1)=2n-2=n+(n-2)>n$ since $n>4$. But clearly $n>n$ is impossible.
\end{answer}
\end{ques}

\begin{ques}
In the second part of the previous proof, why could we say that $a>2$?
\answerline
\begin{answer}
We assumed that $n=a^2>4$, so clearly $a>2$.
\end{answer}
\end{ques}


\begin{exa}\label{exa:AMGM_inequalityn=2}
Prove the
{\em Arithmetic Mean-Geometric Mean Inequality}, which states that for
all non-negative real numbers $x$ and $y$,  
$$\sqrt{xy} \leq \dfrac{x+y}{2}.   $$
\begin{pf}
Since $x$ and $y$ are non-negative, $\sqrt{x}$ and $\sqrt{y}$ are real numbers, so
 $\sqrt{x} - \sqrt{y}$ is a real number.
 Since the square of any real number is greater than
or equal to $0$ we have
$$ (\sqrt{x} - \sqrt{y})^2\geq 0.$$
Expanding (recall the FOIL method?) we get
$$ x - 2\sqrt{xy} + y \geq 0.$$
Adding $ 2\sqrt{xy}$ to both sides and dividing by 2, we get
$$\dfrac{x + y}{2} \geq \sqrt{xy}, $$
yielding the result.
\end{pf}
\end{exa}

The previous example illustrates the creative part of writing proofs.  
The proof started out considering $\sqrt{x} - \sqrt{y}$, which doesn't seem to be related
to what we wanted to prove.  But hopefully after you read the entire proof you see why
it makes sense.  If you are saying to yourself ``I would never have thought of  starting with $\sqrt{x}-\sqrt{y}$?,''
or ``How do you know where to start?,'' I am afraid there are no easy answers.   Writing proofs is as much
of an art as it is a science.  There are three things that can help, though.  
First, {\em don't be afraid to {\bf\em experiment}}.
If you aren't sure where to begin, try starting at the end.  Think about the end goal and work backwards until
you see a connection.  Sometimes working both backward and forward can help.  
Try some algebra and see where it gets you.  But in the end, make sure your proof goes
from beginning to end.  In other words, the order that you figured things out should not necessarily
dictate the order they appear in your proof.

The second thing you can do is to {\bf\em read example proofs}.  
Although there is some creativity necessary in proof writing, 
it is important to follow proper proof writing techniques.  
Although there are often many ways to prove the same statement, 
there is often one technique that works best for a given type of problem. 
As you read more proofs, you will begin to have a better understanding of the various techniques used,
know when a particular technique might be the best choice, 
and become better at writing your own proofs.  If you see several proofs of similar problems,
and the proofs look very similar, then when you prove a similar problem, your proof should probably resemble
those proofs.  This is one area where some students struggle---they submit proofs that look
nothing like any of the examples they have seen, and they are often incorrect.  
Perhaps it is because they are afraid
that they are plagiarizing if they mimic another proof too closely.  However, mimicking a proof is not the
same as plagiarizing a sentence.  To be clear, by `mimic', I don't mean just copy exactly what you see.  I mean
that you should read and understand several examples.  Once you understand the technique used in those examples,
you should be able to see how to use the same technique in your proof.  For instance, in many of the examples
related to even numbers, you may have noticed that they start with statement like {\em ``Assume $x$ is even.  Then
$x=2a$ for some integer $a$.''}  So if you need to write a proof related to even numbers, what sort of 
statement might make sense to begin your proof?

The third thing that can help is {\bf\em practice}.  This applies not only to writing proofs, but to learning
many topics.
An analogy might help here. Learning is often like sports---you don't learn how to play basketball 
(or insert your favorite sport, video game, or other hobby that takes some skill) 
by reading books and/or watching people play it. 
Those things can be helpful (and in some cases necessary), 
but you will never become a proficient basketball player unless you practice. 
Practicing a sport involves running many drills to work on the fundamentals and then applying the 
skills you learned to new situations. Learning many topics is exactly the same.  First you need to 
do lots of exercises to practice the fundamental skills. Then you can apply those skills to new situations. 
When you can do that well, you know you have a good understanding of the topic.  So if you want to become
better at writing proofs, you need to write more proofs.

\begin{ques}
What three things can help you learn to write proofs?
\begin{enumerate}
\item \blanklinefill
\item \blanklinefill
\item \blanklinefill
\end{enumerate}
\begin{answer}
\begin{enumerate}
\item {\bf Experiment}.  If you aren't sure what to do, don't be afraid to try things.
\item {\bf Read Examples}. But don't just read.  Make sure you {\em understand} them.
\item {\bf Practice}.  It makes perfect!  
\end{enumerate}
\end{answer}
\end{ques}

\newpage
%%%%%%%%%%%%%%%%%%%%%%%%%%%%%%%%%%%%%%%%%%%%%%%%%%%%%%%%%%%%%%%%%%%%%%%%%%%%%%%%%%%%%%
\section{Implication and Its Friends}\label{section:implication}
This section is devoted to developing some of the concepts that will be necessary for
us to discuss the ideas behind the next few proof techniques.

%\begin{df}
%\index{proposition}
%\index{boolean!proposition}
%\index{truth value}
%A {\bf boolean proposition} (or simply {\bf proposition}) is a statement which is either {\bf true} or {\bf false}.
%We call this the {\bf truth value} of the proposition.
%\end{df}

Although not technically interchangeable, you may sometimes see the word {\em statement} instead of 
{\em proposition}.  Context should help you determine whether or not a given usage of the word {\em statement}
should be understood to mean {\em proposition}.
We saw the following in the previous chapter, but
it is worth giving the definition again
(stated slightly differently here so it better fits
with our usage in this context). 

\begin{df}
\index{implication}
An {\bf implication} is a proposition of the form ``if $p$, then $q$,''
where $p$ and $q$ are propositions.
$p$ is called the {\bf premise} and $q$ is called the {\bf conclusion}.

It is sometimes written as $p\rightarrow q$, which is read ``$p$ implies $q$.''
It is a statement that asserts that if $p$ is a true proposition then $q$ is a true proposition.

An implication is true unless $p$ is true and $q$ is false.
\end{df}


\begin{exa}
The proposition ``If I do well in this course, then I can take the next course'' is an implication.
However, the proposition ``I can do well in this course and take the next course'' is {\em not} an implication.
\end{exa}

\begin{exa}
Consider the implication 
\begin{quote}
``If you read \xkcd{}, then you will laugh.''\footnote{If you are unfamiliar with \xkcd{}, go to 
\href{http://xkcd.com}{http://xkcd.com}, 
but don't get distracted for too long!}
\end{quote}
If you read \xkcd{} and laugh, you are being consistent with the proposition.
If you read \xkcd{} and {\em do not laugh}, then you are demonstrating that the proposition is false.

But what if you don't read \xkcd{}?  Are you demonstrating that the proposition is true or false?
Does it matter whether or not you laugh?  
It turns out that you are {\em not} disproving it in this case--in other
words, the proposition is still true if you don't read \xkcd{}, whether or not you laugh.  Why?  
Because the statement is not saying anything about laughing by itself. 
It is only asserting that {\bf\em IF} you read \xkcd{}, then you will laugh.
In other words, it is a {\em conditional statement}, with the condition being that 
you read \xkcd{}.  The statement is saying nothing about anything if you don't read \xkcd{}.

So the bottom line is that if you do not read \xkcd{}, the statement is still true.
\end{exa}
\begin{ques}
When is the implication 
``If you read \xkcd{}, then you will laugh'' false?\\
\answerlinetwo
\begin{answer}
Only when you read \xkcd{} and you don't laugh.
\end{answer}
\end{ques}

\begin{exer}
Consider the implication ``If you build it, they will come.''
What are all of the possible ways this proposition could be false?\\
\sollinethree
\begin{answer}
If you build it and they don't come, the proposition is false.  This is the only case where it is false.
To see this, notice that if you build it and they do come, it is true.  If you don't build it, then it
doesn't matter whether or not they come--it is true.
\end{answer}
\end{exer}

Given an implication $p\rightarrow q$, 
there are three related propositions. 
We will introduce each and discuss how they
are related to each other.



%But first we need to recall the negation of a proposition.
%
%\begin{df}
%\index{negation}
%\index{not}
%Given a proposition $p$, the {\bf negation} of $p$, written $\neg p$, is the proposition ``not $p$'' or
%``it is not the case that $p$.''
%\end{df}
%\begin{exa}
%If $p$ is the proposition ``$x \leq y$'' then $\neg p$ is the proposition ``it is not the case that $x\leq y$,'' or
%``$x>y$''.
%\end{exa}

%\begin{rem}
%It is easy to incorrectly negate sentences, 
%especially when they contain words like ``and'', ``or'',
% ``implies'', and ``if.''  This will become easier after we study logic in
%Chapter~\ref{chapter:logicEtc}.
%\end{rem}

\begin{df}
\index{contrapositive}
The {\bf contrapositive} of a proposition of the form ``if $p$, then $q$'' is the proposition  
``if $q$ is not true, then $p$ is not true''  or ``if not $q$, then not $p$''
or $\neg q \rightarrow \neg p$.
\end{df}

\begin{ques}
What is the contrapositive of the proposition ``If you know Java, then you know a programming language''?\\
\answerlinetwo
\begin{answer}
If you don't know a programming language, then you don't know Java.
\end{answer}
\end{ques}

\begin{thm}\label{impConv}
An implication is true if and only if its contrapositive is true.
Stated another way, {\bf an implication and its contrapositive are equivalent}.
\end{thm}

\begin{fib}
Prove Theorem~\ref{impConv}.
\begin{pf}
Let $p\rightarrow q$ be our implication.  According to the definition of implication, 
it is false when $p$ is true and $q$ is false
and \blankline{.75} otherwise.  Put another way, it is true unless $p$ is true and $q$ is false.
The contrapositive, $\neg q\rightarrow \neg p$, is false when 
$\neg q$ is true and \blankline{.5} is false, and true otherwise.  Notice that this is equivalent
to $q$ being \blankline{.5} and \blankline{.5} being true.  
Thus, the contrapositive is true unless \blankline{1.75} and \blankline{1.75}.
But this is exactly when $p\rightarrow q$ is true.
\end{pf}
\begin{answer}
true; $\neg p$; false; $p$; $p$ is true; $q$ is false (the last two can be in either order).
\end{answer}
\end{fib}

\begin{df}
\index{inverse}
The {\bf inverse} of a proposition of the form ``if $p$, then $q$'' is the proposition  
``if $p$ is not true, then $q$ is not true''  or ``if not $p$, then not $q$''
or $\neg p \rightarrow \neg q$.
\end{df}

\begin{ques}\label{ques:javaProp}
What is the inverse of the proposition ``If you know Java, then you know a programming language''?\\
\answerline
\begin{answer}
If you don't know Java, then you don't know a programming language.
\end{answer}
\end{ques}

\begin{ques}
Are a proposition and its inverse equivalent?  Explain, using the proposition from 
Question~\ref{ques:javaProp} as an example.\\
\answerlinethree
\begin{answer}
They are {\em not} equivalent.  Since Java is a programming language, the proposition seems obviously true.
However, what if someone knows C++ but not Java?  Then they know a programming language but they don't
know Java. Thus, the inverse is false.  Since one is true and the other is false, 
the proposition and its inverse are clearly not equivalent.
\end{answer}
\end{ques}

\begin{df}
\index{converse}
The {\bf converse} of a proposition of the form ``if $p$, then $q$'' is the proposition ``if $q$, then $p$''
or $q \rightarrow  p$.
\end{df}

\begin{ques}
What is the converse of the proposition ``If you know Java, then you know a programming language''?\\
\answerline
\begin{answer}
If you know a programming language, then you know Java.
\end{answer}
\end{ques}

\begin{ques}
Are a proposition and its converse equivalent?  Explain using the proposition about Java/programming languages.\\
\answerlinethree
\begin{answer}
They are {\em not} equivalent.  Since Java is a programming language, the proposition seems obviously true.
However, what if someone knows C++ but not Java?  Then they know a programming language but they don't
know Java. Thus, the converse is false.  Since one is true and the other is false, 
the proposition and its converse are clearly not equivalent.
\end{answer}
\end{ques}

As you have just seen, the {\em inverse} and {\em converse} of an implication are not equivalent to
the implication.  However, it turns out that 
the {\em inverse} and {\em converse} of a proposition are equivalent to each other.

\begin{thm}\label{invConv}
The inverse of an implication is true if and only if the converse of the implication is true.
Stated another way, 
{\bf the inverse and converse of an implication are equivalent
to each other}.
\end{thm}

You will be asked to prove Theorem~\ref{invConv} in Problem~\ref{pro:invConc}.  If you think about it in the right way,
it should be fairly easy to prove.
Table~\ref{table:implAndFriends} summarizes what we know.

\begin{table}[tbh]
\centering
$
\begin{array}{|l|l|r|}
\hline
\mbox{Implication	}&p\rightarrow q&\multirow{2}{*}{$\left\}\begin{array}{p{.001in}}\\ \\ \end{array}\right.$Equivalent}\\
\mbox{Contrapositive}&\neg q \rightarrow \neg p&\\
\hline
\mbox{Converse}&q\rightarrow p&\multirow{2}{*}{$\left\}\begin{array}{p{.001in}}\\ \\ \end{array}\right.$Equivalent}\\
\mbox{Inverse}&\neg p\rightarrow\neg q&\\
\hline
\end{array}
$
\caption{Implication and friends}\label{table:implAndFriends}
\end{table}


\begin{exa}
Here is an example of an implication and its friends:
\begin{enumerate}
\item {\bf Implication} If I get to watch ``The Army of Darkness,'' then I will be happy.
\item{\bf Inverse} If I do not get to watch ``The Army of Darkness,'' then I will not be happy.
\item{\bf Converse} If I am happy, then I got to watch ``The Army of Darkness.''
\item{\bf Contrapositive} If I am not happy, then I didn't get to watch ``The Army of Darkness.''
\end{enumerate}
\end{exa}

\begin{ques}
Using the propositions from the previous example, answer the following questions.
\begin{enumerate}[(a)]
\item 
Give an explanation of why an {\em implication} might be true, but the {\em inverse} false.\\
\answerlinethree
\item
Explain why an {\em implication} is saying the exact same thing as its {\em contrapositive}.
(Don't just say ``By Theorem~\ref{impConv}.'')\\
\answerlinethree
\end{enumerate}
\begin{answer}
(a)
The implication states that if I get to watch ``The Army of Darkness'' that I will be happy.  However,
it doesn't say that it is the only thing that will make me happy.  
For instance, if I get to see ``Iron Man'' instead, that would also make me happy.  
Thus, the inverse statement is false.\\
(b)
I will use fact that $p\rightarrow q$ is true unless $p$ is true and $q$ is false.
The implication is true unless I watch ``The Army of Darkness'' and I am not happy.
The contrapositive is ``If I am not happy, then I didn't get to watch `The Army of Darkness.' ''
This is true unless I am not happy and I watched ``The Army of Darkness.''
Since this is exactly the same cases in which the implication are true, the implication
and its contrapositive are equivalent.
\end{answer}
\end{ques}


Implications can be tricky to fully grasp and it is easy to get your head turned around when dealing with them.
We will discuss them in quite a bit of detail throughout the next few sections in order to help
you understand them better.

%%%%%%%%%%%%%%%%%%%%%%%%%%%%%%%%%%%%%%%%%%%%%%%%%%%%%%%%%%%%%%%%%%%%%%%%%%%%%%%%%%%%%%
\newpage
\section{Proof by Contradiction}\label{section:contradiction}

In this section we will see examples of {\em proof by contradiction}.
\index{contradiction proof}
\index{proof!by contradiction}
For this technique, when trying to prove a statement, we assume that its negation
is true and deduce incompatible statements from this (i.e. we prove something we know to be false).  
This implies that the original statement must be true.  

When applied to a proposition, we assume that the premise is true but the conclusion is false,
with the goal still to show that something that is false is true. Since we ``obtained a contradiction,''
it must be that the conclusion is true. We will explain in more detail the idea behind this proof technique
after a few examples.

\begin{exa}
Prove that if $5n+2$ is odd, then $n$ is odd.
\begin{pf}
Assume that $5n+2$ is odd, but that $n$ is even.  Then $n=2k$ for some integer $k$.
This implies that
$5n+2=5(2k)+2 = 10k+2=2(5k+1)$, which is even.  But this contradicts our assumption that
$5n+2$ is odd.  Therefore it must be the case that $n$ is odd.
\end{pf}
The idea behind this proof is that if we are given the fact that $5n+2$ is odd, 
we are asserting that $n$ must be odd.
How do we prove that $n$ is odd?  We could try a direct proof, but it is actually easier to 
use a proof by contradiction in this case.  
The idea is to consider what would happen if $n$ is {\em not} odd.  
What we showed was that if $n$ is not odd, then $5n+2$ has to be even.  
But we {\em know} that $5n+2$ is odd because that was our initial assumption. 
How can $5n+2$ be both odd and even?  It can't. 
In other words, our proof lead to a contradiction--an impossibility.  
Therefore, something is wrong with the proof.  But what?  
If $n$ is indeed even, our proof that $5n+2$ is even is correct.  So there is only
one possible problem--$n$ must not be even.  The only alternative is that $n$ is odd.
Can you see how this proves the statement ``if $5n+2$ is odd, then $n$ is odd?''
\end{exa}

\begin{rem}
If you are somewhat confused at this point that's probably O.K.  Keep reading, and re-read
this section a few times if necessary.  At some point you will have an ``Aha'' moment and
the idea of contradiction proofs will make sense.
\end{rem}

\begin{exa}
Prove that if $n=ab$, where $a$ and $b$ are positive integers, then either $a\leq \sqrt{n}$ or $b\leq \sqrt{n}$.
\begin{pf}
Let's assume that $n=ab$ but that the statement ``either $a\leq \sqrt{n}$ or $b\leq \sqrt{n}$'' is false.  
Then it must be the case that $a>\sqrt{n}$ and $b>\sqrt{n}$.  But then $a b >\sqrt{n}\sqrt{n}=n$.
But this contradicts the fact that $ab=n$.  Since our assumption that  $a>\sqrt{n}$ and $b>\sqrt{n}$
lead to a contradiction, it must be false.
Therefore it must be the case that either $a\leq \sqrt{n}$ or $b\leq \sqrt{n}$.
\end{pf}
\end{exa}

Sometimes your proofs will not directly contradict an assumption made but instead
will contradict a statement that you otherwise know to be true.  For instance, if you ever conclude
that $0>1$, that is a contradiction. 
The next example illustrates this.

\newpage
\begin{fib}
Show, without using a calculator, that 
$\dis{6 - \sqrt{35} < \frac{1}{10}}$.
\begin{pf}Assume that  $\dis{6 - \sqrt{35} \geq \frac{1}{10}}$.
Then  $\dis{6 - \frac{1}{10}\geq}$\blankline{.75}.  
If we multiple both sides by $10$ and do a little
arithmetic, we can see that $59\geq$\blankline{1} . 
Squaring both sides we obtain \blankline{1.25},
which is clearly \blankline{.9}.

\noindent
Thus it must be the case that $\dis{6 - \sqrt{35} < \frac{1}{10}}$.
\end{pf}
\begin{answer}
$ \sqrt{35}$; $10\sqrt{35}$; $3481 \geq 3500$; {\em nonsense} or {\em false} or {\em a contradiction}.
\end{answer}
\end{fib}

Now that we have seen a few examples, let's discuss contradiction proofs a little more formally.\\
Here is the basic concept of contradiction proofs:  You want to prove that a statement $p$ is true. 
You ``test the waters'' by seeing what happens if $p$  is {\em not} true.
So you assume $p$ is false and use proper proof techniques to arrive at a contradiction.
By ``contradiction'' I mean something that cannot possibly be true.  Since you proved something that
is not true, and you used proper proof techniques, then it must be that your assumption was incorrect.
Therefore the negation of your assumption---which is the original statement you wanted to prove---must
be true.

\begin{eval}\label{eval:abeven}
Use the definition of even and odd to prove that if $a$ and $b$ are integers and $ab$ is even, then
at least one of $a$ or $b$ is even.
\begin{spfenum}
\item
By definition of even numbers, let $a$ be an even integer $2n$, and by
the definition of odd numbers, let $b$ be an odd integer $2q+1$.
Then $(2n)(2q+1)=4nq+2n=2(2nq+1)$.
Since $2nq+1$ is an integer, $2(2nq+1)$ is an even integer by definition of even.\\
\evallinethree
\item
If true, either one is odd and the other even, or they are both even, so we will show that
the product of an even and an odd is even, and that the product of two evens integers is even.\\
Let $a=2k$ and $b=2x+1$. 
$(2k)(2x+1)=4kx+2k=2(2kx+k)$.  $2kx+k$ is an integer so $2(2kx+k)$ is even.\\
Let $a=2k$ and $b=2x$.  $(2k)(2x)=4kx=2(2kx)$ since $2kx$ is an integer, $2(2kx)$ is even.\\
Thus, if $a$ and $b$ are integers, $ab$ is even, at least one of $a$ or $b$ is even.\\
\evallinetwo
\item
Let $a$ and $b$ be integers and assume that $ab$ is even, but that neither $a$ nor $b$ is even.
Then both $a$ and $b$ are odd, so $a=2n+1$ and $b=2m+1$ for some integers $n$ and $m$.  But then
$ab=(2n+1)(2m+1)=4nm+2n+2m+1=2(2nm+n+m)+1$, which is odd since $2nm+n+m$ is an integer.
This contradicts the fact that $ab$ is even.  Therefore either $a$ or $b$ must be even.\\
\evallinethree
\end{spfenum}
\begin{answer}
\begin{pfevalenum}
\item
Here are my comments on this proof:
\begin{itemize}
\item
It is proving the wrong thing.  This proves that the product of an even number and
an odd number is even.  But it doesn't even do that quite correctly as we will see next.
\item
The first sentence is phrased weird--we are not letting $a$ be even {\em by} the definition of even.
We are {\em using} the definition.
\item
It does not state that $n$ and $q$ need to be integers.
\item
Although it is not incorrect, using $n$ and $q$ is just weird.
It is customary to use adjacent letters, like $n$ and $m$, or $q$ and $r$.
\item
Given the above problems, I would rephrase the first sentence as
{\em `Let $a$ be an even number and $b$ be an odd number.
Then $a=2n$ and $b=2m+1$ for some integers $n$ and $m$.'}
\item
There is an algebra mistake.  The product should be $2(2nq+n)$.
\item
The last sentence is actually perfect (again, except for the fact that it isn't proving the right thing).
\end{itemize}
\item
This proof is incorrect.  It actually proves the {\em converse} of the statement.  (We'll learn more about
converse statements later.)
In other words, it proves that if at least of one of $a$ or $b$ is even, then $ab$ is even.  
This is {\em not} the same thing.
It is a pretty good proof
of the wrong thing, but it can be improved in at least 4 ways.
\begin{itemize}
\item
It defines $a$ and $b$ but never really uses them.
They should be used at the beginning of the algebra steps (i.e. $a\cdot b=\cdots$) to make
it clear that the algebra is related to the product of these two numbers.
\item
It needs to state that $k$ and $x$ are integers.
\item
As above, using $k$ and $x$ is weird (but not wrong).  It would be better to use $k$ and $l$, or $x$ and $y$.
\item
It needs a few words to bring the steps together.  
In particular, sentences should not generally begin with algebra.
\end{itemize}
Taking into account these things, the second part could be rephrased as follows.
\begin{quote}
Let $a=2n$ and $b=2m+1$, where $n$ and $m$ are integers.
Then $ab=(2n)(2m+1)=4nm+2n=2(2nm+n)$, which is even since $2nm+n$ is an integer.
\end{quote}
\item
This proof is correct.
\end{pfevalenum}
\end{answer}
\end{eval}

For some students, the trickiest part of contradiction proofs is what to contradict.  
Sometimes the contradiction is the fact that $p$ is true.  
At other times you arrive at a statement that is clearly false (e.g. $0>1$).
Generally speaking, you should just try a few things (e.g. do some algebra) and see where it leads.
With practice, this gets easier. 
In fact, with enough practice this will probably become one of your favorite techniques.
When a direct proof doesn't seem to be working this is usually the next technique I try.

\begin{exa}
Let $a_1$, $a_2$, $\ldots$, $a_n$ be real numbers.
Prove that at least one of these numbers is greater or equal to
the average of the numbers.
\begin{pf}
The average of the numbers is
$A=(a_1+a_2+ \ldots + a_n)/n$.
Assume that none of these numbers is greater than or equal to $A$.  That is, 
$a_i<A$ for all $i=1,2,\ldots n$.  Thus
$(a_1+a_2+ \ldots + a_n)<n A$. Solving for $A$, we get
$A > (a_1+a_2+ \ldots + a_n)/n=A$, which is a contradiction.
Therefore at least one of the numbers is greater than or equal to the average.
\end{pf}
\end{exa}



Our next contradiction proof involves {\em permutations}.  Here is the definition and an example
in case you haven't seen these before.

\begin{df}
\index{permutation}
A {\bf permutation} is a function from a finite set to itself that
reorders the elements of the set. 
Since we haven't formally discussed functions yet, the
following informal definition will probably make more sense
to some of you:
a {\bf permutation} of a set of objects is an
ordering of those objects.
\end{df}

\vspace*{-.1in}

\begin{exa}
Let $S$ be the set $\{a, b, c\}$.
Then $(a, b, c)$, $(b, c, a)$ and $(a,c,b)$ are permutations of $S$.
$(a, a, c)$ is not a permutation of $S$ because it repeats $a$ and does not contain $b$.
$(b, d, a)$ is not permutations of $S$ because $d$ is not in $S$, and $c$ is missing.
\end{exa}

\vspace*{-.1in}

\begin{exer}
List all of the permutations of the set $\{1, 2, 3\}$. (Hint:  There are 6.)\\
\answerlinetwo
\begin{answer}
$(1,2,3)$, $(1,3,2)$, $(2,1,3)$, $(2,3,1)$, $(3,1,2)$, and $(3,2,1)$.
\end{answer}
\end{exer}

\begin{rem}
In many contexts, when a list of objects occurs in {\bf curly braces},
the order they are listed does not matter (e.g. $\{a, b, c\}$ and $\{b, c, a\}$ mean the same thing).  
On the other hand, when a list occurs in
{\bf parentheses}, the order {\bf does} matter.
Thus, $(a, b, c)$ and $(b, c, a)$ {\bf do not} mean the same thing. 
\end{rem}

\begin{exa}
Let $(a_1, a_2, \ldots, a_n)$ be an arbitrary
permutation of the numbers $1, 2, \ldots , n, $ where $n$ is an odd number. Prove
that the product $(a_1 - 1)(a_2 - 2)\cdots (a_n - n)$ is even.
\begin{pf}
Assume that the product is odd. Then all of the differences $a_k -
k$ must be odd. Now consider the sum
$S = (a_1 - 1) + (a_2 - 2) + \cdots + (a_n - n).$ 
Since the $a_k$'s are a just a 
reordering of $1, 2, \ldots , n$, $S=0$.
But $S$ is the sum of an odd  number of odd integers, so it must be odd. Since $0$ is not odd, we have
a contradiction.  Thus our initial assumption that all of the $a_k - k$ are odd
is wrong, so at least one of them is even and hence the product is even.
\end{pf}
\end{exa}

\begin{ques}
Why did the previous proof begin by assuming that the product was odd?\\
\answerlinethree
\begin{answer}
Since it wasn't obvious how to do a direct proof of the fact, proof by contradiction seemed like
the way to go.  So we begin by assuming what we want to prove (that the product is even) is false.
The short answer: {\em Because contradiction proofs generally begin by assuming the negation of what
you want to prove.}
\end{answer}
\end{ques}

\begin{ques}
In the previous proof, we asserted that $S=0$.  Why was this the case?\\
\answerlinethree
\begin{answer}
The proof gives the justification for this, but you may have to think about it
for it to entirely sink in.  
Consider carefully the definition of $S$:
$S = (a_1 - 1) + (a_2 - 2) + \cdots + (a_n - n).$ 
Notice it adds and subtracts terms.  If $S=0$, then the amount added and subtracted
must be the same.  And if you think about it for a few minutes, especially in light
of the justification given in the proof, you should see why.
If you can't see it right away, go back to how the $a_k$'s are defined and think a little more.
If you get totally stuck, try an example with $n=3$ or $4$.
\end{answer}
\end{ques}

We will use facts about rational/irrational numbers to demonstrate some of the proof techniques.
In case you have forgotten, here are the definitions.

\begin{df}
\index{rational number}
\index{irrational number}
Recall that 
\begin{itemize}
\item A {\bf rational number} is one that can be written as $p/q$, where $p$ and $q$ are integers,
with $q\neq 0$.
\item An {\bf irrational number} is a real number that is not rational.
\end{itemize}
\end{df}

\begin{exa}\label{sqrt2irr}
Prove that $\sqrt{2}$ is irrational.
We present two slightly different proofs.  In both, we will use the fact that any positive
integer greater than 1 can be factored uniquely as the product of primes (up to the order of the factors).
\begin{pfenum}
\item 
Assume that $\dis{\sqrt{2} = \frac{a}{b}},$ where
$a$ and $b$ are positive integers with $b\neq 0$.  We can assume $a$ and $b$ have no factors in common 
(since if they did, we could cancel them and use the resulting numerator and denominator as $a$ and $b$).
Multiplying by $b$ and squaring both sides yields  $2b^2 = a^2.$ 
Clearly $2$ must be a factor of $a^2$.  Since $2$ is prime,  $a$ must have $2$ as a factor, and 
therefore $a^2$ has $2^2$ as a factor.  Then $2b^2$ must also have $2^2$
as a factor.  But this implies that $2$ is a factor of $b^2$, and therefore a factor of $b$.  
This contradicts the fact that $a$ and $b$ 
have no factors in common.  Therefore $\sqrt{2}$ must be irrational.
\item
Assume that $\dis{\sqrt{2} = \frac{a}{b}},$ where
$a$ and $b$ are positive integers with $b\neq 0$. 
Multiplying by $b$ and squaring both sides yields $2b^2 = a^2.$ Now both $a^2$ and $b^2$ have an
even number of prime factors. So $2b^2$ has an odd number of
primes in its factorization and $a^2$ has an even number of primes
in its factorization. This is a contradiction since they are the same number. 
Therefore $\sqrt{2}$ must be irrational.
\end{pfenum}
\end{exa}

\begin{ques}
In proof 2 from the previous example, why do $a^2$ and $b^2$ have an even number of factors?\\
\answerlinethree
\begin{answer}
Because $a^2=a\cdot a$, so to list the factors of $a^2$ you can list the factors of $a$ twice.
Thus, $a^2$ has twice as many factors as $a$, so it must be an even number.
\end{answer}
\end{ques}

Now that you have seen a few more examples, it is time to begin the discussion about 
how/why contradiction proofs work. 
We will start with the following idea that you may not have thought of before.
It turns out that if you start with a false assumption, then
you can prove that {\em anything} is true.  
It may not be obvious how 
(e.g. How would you prove that all elephants are less than 1 foot tall given that $1+1=1$?), 
but in theory it is possible.
This is because statements of the form ``$p$ implies $q$'' are true if $p$ (called the {\em premise}) is false, 
regardless of whether or not $q$ (called the {\em conclusion}) is true or false.

\begin{exa}
The statement ``If chairs and tables are the same thing, then the moon is made of cheese''
is true.  This may seem weird, but it is correct.  Since chairs and tables are not the same thing,
the premise is false so the statement is true.  But it is important to realize that the fact
that the whole statement is true doesn't tell us anything about whether or not the moon is made of cheese.
All we know is that {\em if} chairs and tables were the same thing, then the moon {\em would have to} 
be made out of cheese in order for the statement to be true.
\end{exa}

\begin{exa}
Consider what happens 
if your parents tell you ``If you clean your room, then we will take you to get ice cream.''
If you don't clean your room and your parents don't take you for ice cream, 
did your parents tell a lie? No.  What if they {\em do} take you for ice cream?  They still haven't
lied because they didn't say they wouldn't take you if you didn't clean your room.
In other words, if the premise is false, the whole statement is true regardless of whether or not
the conclusion is true.
\end{exa}

It is important to understand that when we say that a statement of the form ``$p$ implies $q$'' is true, 
we are {\em not} saying that $q$ is true.  We are only saying that {\em if $p$ is true, then $q$ has to be true}.  
We aren't saying anything about $q$ by itself.  
So, if we know that ``$p$ implies $q$'' is true, and we also know that $p$ is true, then $q$ must also
be true.  This is a rule called {\em modus ponens}\index{modus ponens},
and it is at the heart of contradiction proofs as we will see shortly.

\begin{exer}
It might help to think of statements of the form ``$p$ implies $q$'' as rules 
where breaking them is equivalent to the statement being false.
For instance, consider the statement {\em ``If you drink alcohol, you must be 21.''}
If we let $p$ be the statement ``you drink alcohol'' and $q$ be the statement ``you are 21,''
the original statement is equivalent to ``$p$ implies $q$''.  
\begin{enumerate}
\item If you drink alcohol and you are 21, did you break the rule?
\sblanklinefill
\item If you drink alcohol and you are not 21, did you break the rule?
\sblanklinefill
\item If you do not drink alcohol and you are 21, did you break the rule?
\sblanklinefill
\item If you do not drink alcohol and you are not 21, did you break the rule?
\sblanklinefill
\item
Generalize the idea.  If you have a statement of the form ``$p$ implies $q$'',
where $p$ and $q$ can be either true or false statements, exactly when can the statement be false?\\
\blanklinefill
\item
If you do not drink alcohol, does it matter how old you are?
\sblanklinefill
\item
Can a statement of the form ``$p$ implies $q$'' be false if $p$ is false?  Explain.\\
\blanklinefill\\
\blanklinefill
\end{enumerate}
\begin{answer}
(1) No. (2) Yes. (3) No. (4) No. 
(5) Statements of the form ``$p$ implies $q$'' are false
precisely when $p$ is true and $q$ is false.
(6)  No.  Whether or not you are 21, you aren't breaking the rule.
(7) No.  If $p$ is false, whether or not $q$ is true or false doesn't matter--the statement is true.
Let's consider the previous question--if you do not drink alcohol, you are following the rule
regardless of whether or not the statement ``you are 21'' is true or false.
\end{answer}
\end{exer}

Now we are ready to explain the idea behind contradiction proofs.  
We want to prove some statement $p$ is true. 
We begin by assuming it is false---that is, we assume $\neg p$ is true.
We use this fact to prove that $q$---some false statement---is true.  In other words, we
prove that the statement ``$\neg p$ implies $q$'' is true, where $q$ is some false statement.
But if $\neg p$ is true, and ``$\neg p$ implies $q$'' is true, modus ponens tells us that $q$ is true.
Since we know that $q$ is false, something is wrong.
We only have two choices: either $\neg p$ is false or ``$\neg p$ implies $q$'' is false.
If we used proper proof techniques to establish that ``$\neg p$ implies $q$'' is true, 
then that isn't the problem.  Therefore, $\neg p$ must be false, implying
that $p$ is true.  That is why contradiction proofs work.

Let's analyze the second proof from Example~\ref{sqrt2irr} in light of this discussion.
The {\em only} assumption we made was that $\sqrt{2}$ is rational ($\neg p$=``$\sqrt{2}$ is rational'').
From this (and only this), we were able to show that $a^2$ has both an even and an odd number
of factors ($q$=``$a^2$ has an even and an odd number of factors'', and we showed that ``$\neg p$ implies $q$''
is true).
Thus, we know for certain that if $\sqrt{2}$ is rational, 
then $a^2$ has an even and an odd number of factors.\footnote{We did not prove
that $a^2$ has an even and an odd number of factors.  We proved that {\em if $\sqrt{2}$ is rational}, 
then $a^2$ has an even and an odd number of factors.  It is very important that you understand the
difference between these two statements.}
This fact is indisputable since we proved it.  
If it is also true that $\sqrt{2}$ is rational, modus ponens implies that 
$a^2$ has an even and an odd number of factors.  This is also indisputable.  
But we know that $a^2$ cannot have both an even and odd number of factors.  
In other words, we have a contradiction.  Therefore, something that has been said somewhere is wrong.
Everything we said is indisputable except for one thing--that $\sqrt{2}$ is rational. 
That was never something we proved---we just assumed it. 
So it has to be the case that this is false, which means that $\sqrt{2}$ must be irrational.

To summarize, if you want to prove that a statement is true using a contradiction proof, 
assume the statement is false, use this assumption to get a contradiction (i.e. prove a false statement),
and declare that it must therefore be true.

Notice that what $q$ is doesn't matter.  In other words, given the assumption 
$\neg p$, the goal in a contradiction proof is to establish that {\em any} false statement is true.
This is both a blessing and a curse.  The blessing is that any contradiction will do.  
The curse is that we don't have a clear goal in mind, so it can sometimes be difficult to know what to do.
As mentioned previously, this becomes easier as you read and write more proofs.

If this discussion has been a bit confusing, try re-reading it.  The better you understand the theory
behind contradiction proofs, the better you will be at writing them.  
We will revisit some of these concepts in the chapter on logic, so the more you understand from here, the
better off you will be when you get there.
O.K., enough theory.  Let's see some more examples!

\begin{fib}
Let $a, b$ be real numbers.
Prove that if  $ a < b + \epsilon$ for all $\epsilon > 0$, then $a \leq b.$
\begin{pf}
We will prove this by contradiction.  
Assume that \blankline{.75}.\footnote{Hint: What assumption do we always make when doing a contradiction proof?} 
Subtracting $b$ from both sides and dividing by $2$, we get \blankline{1}$>0$. 
Since the inequality $a < b + \epsilon$ holds for
every $\epsilon > 0$ in particular it holds for $\epsilon=$\blankline{1}.\footnote{Same as the previous blank} 
This implies that 
$$a < b + \frac{a - b}{2}=\blankline{2}.$$
If we \blankline{3} (to the previous equation), we obtain $a < b$.
But we started with the assumption that \blankline{1} which is a \blankline{1}.
Therefore, \blankline{2}.
\end{pf}
\begin{answer}
$a > b$; $\frac{a - b}{2}$; $\frac{a - b}{2}$; $b+\frac{a}{2}-\frac{b}{2}=\frac{a}{2}+\frac{b}{2}$; 
subtract $\frac{a}{2}$ from both sides and multiple both sides by $2$;
$a > b$; contradiction; $a\leq b$.
\end{answer}
\end{fib}

The following beautiful proof goes back to Euclid.
It uses the assumption that any integer greater than 1 is either a prime or a product of primes. 

\begin{exa}[Euclid]
Show that there are infinitely many prime numbers.
\label{exa:inf_primes}
\begin{pf}
Assume that there are only a finite number of primes and 
label the primes 
$\{p_1, p_2, \ldots , p_n\}$. 
Consider the number
    $$N = p_1p_2\cdots p_n + 1.$$ This is a positive integer that is clearly greater than 1. 
    Observe that none of the primes on the list $\{p_1, p_2, \ldots , p_n\}$ divides $N$, 
    since division by any of these primes
    leaves a remainder of 1. 
    Since $N$ is larger than any of the primes on this list, it is either a prime
    or divisible by a prime outside this list. 
    But we assumed the list above contained all of the prime numbers.
    This is a contradiction.  Therefore there must be infinitely many primes.
\end{pf}
\end{exa} 

\begin{fib}\label{fibRatPoly}
If $a, b, c$ are odd integers, prove that $ax^2 + bx + c = 0$ does
not have a rational number solution.
\begin{pf}
Suppose $\dfrac{p}{q}$ is a rational solution
to the equation. We may assume that $p$ and $q$ have no prime
factors in common, so either $p$ and $q$ are both odd, or one is
odd and the other even. 
Since $\dfrac{p}{q}$ is a solution, we know that 
$$\blankline{3} = 0.$$
If we \blankline{2}, we obtain  $ap^2 + bpq + cq^2 = 0.$

If both $p$ and $q$ are odd, then $ap^2 + bpq + cq^2 $ is 
\blankline{1} which contradicts the fact that it is \blankline{.5}.

If $p$ is even and  $q$ odd, then  \blanklinefill\\ \blanklinefill.

If $p$ is odd and  $q$ even, then  \blanklinefill\\ \blanklinefill.

Since all possibilities leads to a contradiction, \blanklinefill.

\blanklinefill\\
\end{pf}
\begin{answer}
$a\left(\frac{p}{q}\right)^2 + b\left(\frac{p}{q}\right) + c$;
multiple both sides by $q^2$; 
odd;
0;
$ap^2 + bpq$ is even and $cq^2$ is odd, so  $ap^2 + bpq + cq^2$ is odd;
$bpq + cq^2$ is even and $ap^2$ is odd, so  $ap^2 + bpq + cq^2$ is odd;
$ax^2 + bx + c = 0$ does not have a rational solution if $a$, $b$, and $c$ are odd.
\end{answer}
\end{fib}

One final note on contradiction proofs: 
Only use one when you really need it.
If a direct proof will work, use it. If you use a contradiction proof
instead, you will just be making the proof more complicated
for no good reason. Some students seem to grab onto 
contradiction proofs and try to use it for everything, but they are
not the best choice in many cases. 

\newpage
\section{Proof by Contraposition}\label{section:contraposition}

Consider the statement ``If it rains, then the ground will get wet.''  
It should be pretty easy to convince yourself that this is essentially equivalent to the statement 
``If the ground is not wet, then it didn't rain.''
In fact, since the second statement is just the contrapositive of the first, 
Theorem~\ref{impConv} tells us that they are equivalent.
Again, by {\em equivalent} we simply mean that
either both statements are true or both statements are false.  
This is the idea behind the proof technique in this section.

\begin{df}
\index{proof!contrapositive}
\index{contraposition!proof by}
A {\bf proof by contraposition} is a proof of a statement of the form ``if $p$, then $q$''  that proves
contrapositive statement instead.  
That is, it proves the equivalent statement ``if not $q$, then not $p$.''
\end{df}

\begin{exa}
Prove that if $5n+2$ is odd, then $n$ is odd.
\begin{pf}
We will instead prove that if $n$ is even (not odd), then $5n+2$ is even (not odd).  Since
this is the contrapositive of the original statement,  a proof of this will prove that that
the original statement is true.

Assume $n$ is even.  The $n=2a$ for some integer $a$.  Then $5n+2=5(2a)+2 = 2(5a+1)$.
Since $5a+1$ is an integer, $2(5a+1)$ is even.
\end{pf}
\end{exa}

Be careful with proof by contraposition. Do not make the mistake of trying to
prove the {\em converse} or {\em inverse} instead of the {\em contrapositive}.
In that case, you may (sometimes) write a correct proof, but it would be a proof of the wrong thing.

In the next example we will see the similarities and differences between contradiction
proofs and proofs by contraposition.

\begin{exa}
Prove that if $5n+2$ is even, then $n$ is even.

\bigskip 

\noindent
\begin{tabular}{lll}
\begin{minipage}{2.65in}
\noindent
{\bf Proof by contraposition:}

We will prove the equivalent statement that if $n$ is odd, then $5n+2$ is odd.

Assume $n$ is odd.  Then $n=2k+1$ for some integer $k$. Then we have that
\begin{eqnarray*}
5n+2&=&5(2k+1)+2\\
&=&10k+5+2\\
&=&10k+7\\
&=&2(5k+3)+1
\end{eqnarray*}
Since $5k+3$ is an integer, this shows that $5n+2$ is odd.

\bigskip

\bigskip

\end{minipage}
&\, \,
&
\begin{minipage}{2.65in}
{\bf Proof by contradiction:}

Assume that $5n+2$ is even but that $n$ is odd.
Since $n$ is odd, $n=2k+1$ for some integer $k$.
Therefore 
\begin{eqnarray*}
5n+2&=&5(2k+1)+2\\
&=&10k+5+2\\
&=&10k+7\\
&=&2(5k+3)+1
\end{eqnarray*}
which is odd since $5k+3$ is an integer.  
But we assumed that $5n+2$ was even, which is a contradiction.
Therefore our assumption that $n$ is odd must be incorrect, so $n$ is even.
\end{minipage}
\end{tabular}
\end{exa}

\begin{eval}
Let $n$ be an integer. 
Use the definition of even/odd to 
prove that if $3n+2$ is even, then $n$ is even using a proof by contraposition.
\begin{spfenum}
\item
We need to show that if $n$ is even, then $3n+2$ is even.
If $n$ is even, then $n=2k$ for some integer $k$.  
Then $3n+2= 3(2k+2) = 6k+6=2(3k)+2(3)$, which is even because it is the sum of two even integers. \\
\evallinethree
\item
We need to show that if $n$ is odd, then $3n+2$ is odd.  
If $n$ is odd then $n=2k+1$ for some integer $k$.
Then $3n+2=3(2k+1)+2
=6k+3+2
=6k+5
=5(\frac{6}{5}k+1)$, which is clearly odd. \\
\evallinethree
\item
We need to show that if $n$ is odd, then $3n+2$ is odd.  If $n$ is odd then $n=2k+1$ for some integer $k$.
Then $3n+2=3(2k+1)+2
=6k+5$,
which is odd by the definition of odd. \\
\evallinethree
\end{spfenum}
\begin{answer}
\begin{pfevalenum}
\item
This is attempting to prove the {\em converse}, not the contrapositive.
Since the converse of a statement is not equivalent to the original statement, this is not a valid proof.
Further, the proof contains an algebra mistake.
Finally, it uses the property that the sum of two even integers is even.  Although this is true, the problem
specifically asked to prove it using the definition of even/odd.
\item
This proof starts out correctly by using the contrapositive statement and the definition of odd.
Unfortunately, the writer claims that $\displaystyle 5(\frac{6}{5}k+1)$ is `clearly odd.'
This is not at all clear.  What about this number makes it odd?  Is it expressed as $2a+1$ for
some integer $a$?  No.  
Even worse, there is a fraction in it, obscuring the fact that the number is even an integer.
\item
This proof is {\em really close}.  The only problem is that we don't know that $6k+5$ is odd 
{\em using the definition of odd}.  All the writer needed to do is take their algebra a little
further to obtain $2(3k+2)+1$, which is odd by the definition of odd since $3k+2$ is an integer.
\end{pfevalenum}
\end{answer}
\end{eval}

\newpage
\section{Other Proof Techniques}\label{section:other}
There are many other proof techniques.  
We conclude this chapter with a small sampling of the more common and/or interesting ones.  
We will see a few other important proof techniques later in the book.

\begin{df}
\index{proof!trivial}
\index{trivial proof}
A {\bf trivial proof} is a proof of a statement of the form ``if $p$, then $q$''  that doesn't use $p$ in the proof.
\end{df}

\begin{exa}
Prove that if $x>0$, then $(x+1)^2 - 2x > x^2$.
\begin{pf}
It is easy to see that 
\begin{eqnarray*}
(x+1)^2 - 2x&=&(x^2+2x+1) - 2x\\
&=&x^2+1\\
&>&x^2.
\end{eqnarray*}
Notice that we never used the fact that $x>0$ in the proof.
\end{pf}
\end{exa}

\begin{df}
\index{counterexample!proof by}
\index{proof!by counterexample}
A {\bf proof by counterexample} is used to disprove a statement by giving an example of it being false.
\end{df}

\begin{exa}
Prove or disprove that the product of two irrational numbers is irrational.
\begin{pf}
We showed in Example \ref{sqrt2irr} that $\sqrt{2}$ is irrational.  But $\sqrt{2}*\sqrt{2}=2$, which
is an integer so it is clearly rational.  Thus the product of 2 irrational number is not always irrational.
\end{pf}
\end{exa}

\begin{exa}
Prove or disprove that ``Everybody Loves Raymond'' (or that ``Everybody Hates Chris'').
\begin{pf}
Since I don't really love Raymond (I also don't hate Chris, in case you care), the statement is clearly false.
\end{pf}
\end{exa}

\begin{exer}
Prove or disprove that the sum of any two primes is also prime.\\
\pflinethree
\begin{answer}
Answers will vary greatly, but one proof is: $3$ and $5$ are prime but $3+5=8=2^3$ is clearly not prime.
\end{answer}
\end{exer}

\begin{df}
\index{proof!by cases}
A {\bf proof by cases} breaks up a statement into multiple cases and proves each one separately.
\end{df}

We have already seen several examples of proof by cases (e.g. Examples~\ref{ndivfactExer} and \ref{fibRatPoly}), but
it never hurts to see another example.
\begin{exa}
Prove that if $x\neq 0$ is a real number, then $x^2>0$.
\begin{pf}
If $x\neq 0$, then either  $x>0$ or $x<0$.

If $x>0$ (case 1), then we can multiply both sides of $x>0$ by $x$, giving $x^2>0$.

If $x<0$ (case 2), then we can write y=-x, where $y>0$.  Then $x^2=(-y)^2=(-1)^2y^2=y^2>0$
by case 1 (since $y>0$).  Thus $x^2>0$.
In either case, we have shown that $x^2>0$.
\end{pf}
\end{exa}

\begin{fib}
Let $s$ be a positive integer. Prove that the closed interval
$[s,2s]$ contains a power of $2$.
\begin{pf}
If $s$ is a power of 2 then \blanklinefill

If $s$ is not a power of 2, then it is strictly between two powers of 2.
That is, $2^{r-1} < s < 2^r$ for some integer $r$.  Then \blanklinefill

\blanklinefill

\blanklinefill

\end{pf}
\begin{answer}  $2s$ is a power of two that is in the closed interval.;
$2^r=2\cdot 2^{r-1}< 2s<2\cdot 2^r=2^{r+1}$, so 
$s< 2^r < 2s < 2^{r+1}$, and so the interval $[s,2s]$ contains $2^r$, 
a power of $2$.
\end{answer}
\end{fib}

\newpage
\section{If and Only If Proofs}\label{section:iff}

Sometimes we will run into ``if and only if''  (abbreviated {\em iff}) statements.
That is, statements of the form  {\bf  $p$ if and only if $q$}.
This is equivalent to the statement ``$p$ implies $q$ and $q$ implies $p$.''
Thus, to prove that an iff statement is true, you need to prove a statement and
its {\em converse}.
``$p$ implies $q$'' is sometimes called the {\em forward direction} and the converse
is sometimes called the {\em backwards direction}.
Sometimes the converse statement is proven by {\em contraposition},
so that instead of proving $q$ implies $p$, $\neg p$ implies $\neg q$ is proven.

\begin{ques}
Why is there a choice between proving {\bf $q$ implies $p$} and proving {\bf $\neg p$ implies $\neg q$} when proving the backwards direction?\\
\answerlinetwo
\begin{answer}
Because these statements are contrapositives of each other.  In other words, they are equivalent.
Therefore you can prove either form of the statement.
\end{answer}
\end{ques}

\begin{exa}
Prove that $x$ is even if and only if $x+10$ is even.
\begin{pf}
If $x$ is even, then $x=2k$ for some integer $k$. 
Then $x+10=2k+10=2(k+5)$.  Since $k+5$ is an integer, then 
$x+10$ is even.  
Conversely, if $x+10$ is even, then $x+10=2k$ for some integer $k$.  
Then $x=(x+10)-10=2k-10=2(k-5)$.  Since $k-5$ is an integer, then $x$ is even.
Therefore $x$ is even iff $x+10$ is even.
\end{pf}
\end{exa}

As we have mentioned before, the examples in this section are quite trivial and may seem
ridiculous--since they are so obvious, why are we bothering to prove them?  The point is
to use the proof techniques we are learning.  We will use the techniques on more complicated
problems later.  For now we want the focus to be on proper use of the techniques.  That is
more easily accomplished if you don't have to think too hard about the details of the proof.

\begin{exer}
Prove that $x$ is odd iff $x+20$ is odd using direct proofs for both directions

\vspace*{2.5in}
\begin{answer}
If $x$ is odd, then $x=2k+1$ for some integer $k$.  
Then $x+20=2k+1+20=2(k+10)+1$, which is odd since $k+10$ is an integer.
If $x+20$ is odd, then $x+20=2k+1$ for some integer $k$.  Then $x=(x+20)-20=2k+1-20=2(k-10)+1$,
which is odd since $k-10$ is an integer.
Therefore $x$ is odd iff $x+20$ is odd.
\end{answer}
\end{exer}

\begin{exer}
Prove that $x$ is odd iff $x+20$ is odd using 
using a direct proof for the forward direction and a proof by contraposition for the backward direction.

\vspace*{2.5in}
\begin{answer}
If $x$ is odd, then $x=2k+1$ for some integer $k$.  
Then $x+20=2k+1+20=2(k+10)+1$, which is odd since $k+10$ is an integer.
If $x$ is even, then $x=2k$ for some integer $k$. 
Then $x+20=2k+20=2(k+10)$.  Since $k+10$ is an integer, then 
$x+20$ is even. Therefore $x$ is odd iff $x+20$ is odd. 
\end{answer}
\end{exer}


\begin{fib}
The two most common ways to prove $p$ iff $q$ are
\begin{enumerate}
\item
Prove that \blankline{2} and \blankline{2}, or
\item
Prove that \blankline{2} and \blankline{2}.
\end{enumerate}
\begin{answer} 
$p$ implies $q$; 
$q$ implies $p$;  
$p$ implies $q$; 
$\neg p$ implies $\neg q$
\end{answer}
\end{fib}

\begin{eval}
Use the definition of odd to prove that $x$ is odd if and only if $x-4$ is odd.
\begin{spfenum}
\item
Assume $x$ is odd. Then $x=2k+1$ for some integer $k$.  Then
$x-4=2k+1-4=2k-3$, which is odd.
Now assume that $x-4$ is odd. Since $(2k+1)-4$ is odd, then $x=2k+1$ is clearly odd.\\
\evallinethree
\item
Assume $x$ is odd. Then $x=2k+1$, so $x-4=(2k+1)-4=2(k-2)+1$, which is odd since
$k-2$ is an integer.  
Now assume $x-4$ is even. Then $x-4=2k$ for some integer $k$.  
Then $x=2k+4=2(k+2)$, which is even since $k+2$ is an integer.\\
\evallinetwo
\end{spfenum}
\begin{answer}
\begin{pfevalenum}
\item
For the forward direction, they didn't use the definition of odd.  Otherwise, that part is fine.
For the backward direction, their proof is nonsense.  They {\em assumed} that $x=2k+1$ when they wrote
$(2k+1)-4$ in the second sentence.  This need to be {\em proven}.
\item
For the forward direction, they didn't specify that $k$ was an integer.  Otherwise it is correct.
The second part of the proof is {\em not} proving the converse.  It is proving the forward direction
a second time using a proof by contraposition.  In other words, this proof just proves the forward
direction twice and does not prove the backward direction.
\end{pfevalenum}
\end{answer}
\end{eval}

\newpage
\section{Common Errors in Proofs}\label{section:proofErrors}
If you arrive at the right conclusion, does that mean your proof is correct?  Some students seem to think
so, but this is absolutely false.  Let's consider the following example.
\begin{exa}
Is the following proof that $\displaystyle \frac{16}{64}=\frac{1}{4}$ correct?  Why or why not?
\begin{spf}
This is true because if I cancel the $6$ on the top and the bottom, I get
$\displaystyle \frac{16}{64}=\frac{1\cancel{6}}{\cancel{6}4}=\frac{1}{4}.$
\end{spf}
\begin{quote}
{\bf Evaluation:} You probably know that you can't cancel arbitrary digits in a fraction, so this
is not a valid proof.  In addition, consider this:  
If this proof is correct, then it could be used to prove that
$\displaystyle \frac{16}{61}=\frac{1\cancel{6}}{\cancel{6}1}=\frac{1}{1}=1$, which is clearly false.
\end{quote}
\end{exa}

\begin{rem}
The point of the previous example is this: 
Don't confuse the fact that what you are trying to prove is true with whether or not your proof
actually {\em proves} that it is true.  
An incorrect proof of a correct statement is no proof at all.
\end{rem}

One rookie mistake that I see often is {\em proof by example}, where the writer attempts to prove 
something in general by proving it for one particular case and assuming it must therefore work 
for all of the other cases.

\begin{ques}\label{ques:badone}
 What is wrong with this `proof' that the sum of two even integers is even?
\begin{spf}
Let $x$ and $y$ be even integers.  Assume $x=4$ and $y=6$, which are both even.
Then $x+y=10=2*5$, which is even since $5$ is an integer. 
Thus, the sum of two even integers is even.
\end{spf}
\answerlinetwo
\begin{answer}
This only proves that $4+6$ is even. It says nothing about the sum of any other two even numbers.
\end{answer}
\end{ques}

Just because a proof seems work out, it does not mean that it is a proof of the correct statement. 
For instance, the proof in Question~\ref{ques:badone} is a correct proof of the fact that the sum of 4 and 6 is even. 
But it is certainly {\em not} a proof that the sum of {\em any} two even numbers is even. 

Let's see an example of a supposed proof of something that is not even true.
Hopefully I do not need to convince you that the proof cannot be valid (since the statement is false).

\begin{exa}
 What is wrong with this `proof' that one more than an even number is divisible by 3?
\begin{spf}
Notice that $14+1=15=3*5$ which is clearly divisible by $3$. Since $14=2*7$ is even,
we just showed that one more than an even number is divisible by 3.
\end{spf}
\begin{quote}
{\bf Evaluation:}
This only shows that one more than 14 is divisible by 3. Notice that $10$ is even, but $10+1=11$ is
not divisible by 3, so the statement that is supposedly being proven here is clearly not true!
\end{quote}
\end{exa}

Hopefully this example helps you see the problem with {\em proof by example}. If the technique worked,
then the proof in the previous example is a valid proof of the false statement that one more than an even number is 
divisible by 3. But since that statement is false, it can't have been a valid proof. Indeed, as we
already mentioned, the proof does show that the statement is true for the given even number (in this case, 14),
but that does not imply anything about the validity of the statement for any other even numbers.

If you want to prove something for a general collection of
numbers (e.g. even number, integers, etc.), then your proof has to be general enough to include all 
possible values. For instance, if you want to prove something about odd numbers, then you let 
$x=2k+1$ where $k$ is an integer. Notice that no matter which odd integer you want to consider, you can
pick $k$ to obtain that value. Thus, if you prove something about the value $x=2k+1$, then
you have proven it for all odd values of $x$. However, if you show it is true for $x=7$ (for instance), you have
only shown that it is true for $x=7$.

Another common mistake when writing proofs is to make one or more invalid assumptions without realizing it. 
This is another case where you end up proving a different statement (usually a more specific statement) 
than the one you set out to prove.
The problem is that when you make this sort of mistake, the proof can sometimes seems to ``work'' 
because you get the conclusion you want. 
Thus, your proof might actually be a valid proof, but it is of the wrong statement.
Thus, it isn't always obvious that you even made a mistake. 

The next few examples should illustrate what can go wrong if you aren't careful.

\begin{ques}
 What is wrong with this `proof' that the sum of two even integers is even?
\begin{spf}
Let $x$ and $y$ be even integers.  Then $x=2a$ for some integer $a$ and $y=2a$
for some integer $a$.
So $x+y=2a+2a=2(a+a)$.  Since $a+a$ is an integer, $2(a+a)$
is even, so the sum of two even integers is even.
\end{spf}
\answerlinetwo
\begin{answer}
The problem is that this is actually a proof that $x+x$ is even if $x$ is even since $x=2a=y$ was assumed.
\end{answer}
\end{ques}

Since the statement in the previous example is true, it can be difficult to appreciate why the proof is wrong.
The proof seems to prove the statement but as you saw in the solution, it actually doesn't.  It proves 
a more specific statement (In this case, it is a proof of the fact that the sum of an even number with itself is even
when it was supposed to be a proof of the fact that the sum of any two even numbers is even.).  

If it seems like we are being too nit-picky, consider the next example which gives a supposed proof
that the sum of two even numbers is divisible by 4 (hopefully you can quickly convince yourself that this
is not a true statement).

\newpage

\begin{ques}
What is wrong with the following `proof' that the sum of two even integers is divisible by 4?
\begin{spf}
Let $x$ and $y$ be two even integers.  Then $x=2a$ for some integer $a$ and $y=2a$
for some integer $a$.
So $x+y=2a+2a=4a$.  Since $a$ is an integer, $4a$ is divisible by $4$, so the sum of
two even integers is divisible by $4$.
\end{spf}
\answerlinetwo
\begin{answer}
Notice that $4$ and $6$ are even, but $4+6=10$ is not divisible by $4$.  So clearly the statement
is incorrect.  Therefore, there must be something wrong with the proof.  
The problem is the same as in the previous example--the proof assumed $x=y$, even if that was
not the intent of the writer.  So what was proven was that if $x$ is even, then $x+x$ is divisible by $4$.
\end{answer}
\end{ques}

Another common mistake students make when trying to prove an identity/equation is to start with what
they want to prove and work both sides of it until they demonstrate that they are equal.  I want to stress that
{\em this is an invalid proof technique}.  Again, if this seems like I am making something out of nothing,
consider this example:

\begin{ques}\label{minusOneIsNotOne}
Consider the following supposed proof that $-1=1$.
\begin{spf}
\begin{eqnarray*}
-1&=&1\\
(-1)^2&=&1^2\\
1&=&1
\end{eqnarray*}
Therefore $-1=1$.
\end{spf}
How do you know that this proof is incorrect?
(Think about the obvious reason, not any technical reason.)\\
\answerline
\begin{answer}
Since it should be clear that the result ($-1=1$) is false, the proof can't possibly be correct.
\end{answer}
\end{ques}

Notice that each step of algebra in the previous proof is correct.  For instance, if $a=b$, then $a^2=b^2$
is correct.
And $(-1)^2$ and $1^2$ are both equal to $1$.  So the majority of the proof contains proper techniques.
It contains just one problem: It starts by assuming something that isn't true.
Unfortunately, one error is all it takes for a proof to be incorrect.

\begin{rem}
When writing proofs, {\bf never} assume something that you don't already know to be true!
In particular, if you are trying to prove an equality, never start with the equality and work both sides 
until you get the same thing. As demonstrated in the previous example, this is not a valid proof technique.
\end{rem}

\newpage

\begin{ques}
When you are given an equation to prove, 
should you prove it by writing it down and working both sides
until you get them both to be the same?  Why or why not?\\
\answerlinetwo
\begin{answer}
No!  Example~\ref{minusOneIsNotOne} should have made it clear that this approach is flawed.
\end{answer}
\end{ques}

Let's be clear about this issue.  
If you known an equation is correct, you can work both sides of it until you get to some desired conclusion.
However, if you have an equation and you do not know whether or not it is correct, 
you cannot start your proof by considering that equation.
As Example~\ref{minusOneIsNotOne} demonstrated, if an equation is {\em not correct}, 
sometimes you can work both sides until they are the same, which gives the illusion that you
have proven that it is correct, which is clearly not possible. Hopefully this makes it clear to 
you that beginning a proof with an unknown equation (e.g. the equation you are trying to prove)
and using it in your proof is not valid. 

\begin{ques}
You are given an equation.  You work both sides of it until they are the same.  
Should you now be convinced that the equation is correct? Why or why not?\\
\answerlinetwo
\begin{answer}
No, you should not be convinced.  
As we just mentioned, whether or not the equation is true, sometimes you can work both
sides to get the same thing.  Thus the technique of working both sides is not valid.  
It doesn't guarantee anything unless you already know that the equation is valid.
\end{answer}
\end{ques}

\begin{rem}
If you already know that an equation is true, then working both sides of it 
(for some purpose other than demonstrating it is true) {\em is} a valid technique.
However, it is more common to start with a known equation 
and work just one side until it is what we want.
\end{rem}

There are plenty of other common errors in proofs.  We will see more examples of them throughout
the remainder of the book (although we will focus more on correct proof techniques!), 
especially in the {\em Evaluate} examples.
I want to say that you will likely see other examples of errors in proofs as you write your own proofs,
but that would be mean. Probably accurate, but still mean.

\newpage
\section{More Practice}\label{section:morePractice}
Now you will have a chance to practice what you have learned throughout this chapter with some more
exercises.  Now that they aren't in a particular section, you will have to figure out what technique to use.

\begin{exer}
Let $p<q$ be two {\em consecutive} odd primes (two primes with no other primes between them). 
Prove that $p+q$ is a composite number.  
Further, prove that it has at least three, not necessarily distinct, prime factors.
(Hint: think about the average of $p$ and $q$.)\\
Proof:\\
\vspace*{1.5in}
\begin{answer}
Since $p$ and $q$ are odd, we know that $p+q$ is even, and so
$\dfrac{p+q}{2}$ is an integer. But $p< q$ gives $2p < p+q < 2q $
and so $p < \dfrac{p+q}{2} < q$, that is, the average of $p$ and $q$
lies between them. Since $p$ and $q$ are consecutive primes, any
number between them is composite, and so divisible by at least two
primes. So $p+q = 2\left(\dfrac{p+q}{2}\right)$ is divisible by the
prime $2$ and by at least two other primes dividing
$\dfrac{p+q}{2}$.
\end{answer}
\end{exer}


\begin{eval}
Prove or disprove that if $x$ and $y$ are rational, then $x^y$ is rational.
\begin{spfenum}
\item
Because $x$ and $y$ are both rational, assume $x=a/b$ where $a$ and $b$ are integers and $b\neq 0$.
We can assume that $a$ and $b$ have no factors in common (since if they did we could cancel them and use
the resulting numbers as our new $a$ and $b$).
Then $x^y=\frac{a^y}{b^y}$, so $x^y$ is rational.\\
\evallinethree
\item
Notice that $x^y$ is just $x$ multiplied by itself $y$ times.  A rational number multiplied
by a rational number is rational, so $x^y$ is rational.\\
\evallinethree
\end{spfenum}
\begin{answer}
\begin{pfevalenum}
\item
This is not correct.  It needs to be shown that $x^y$ can be written as $c/d$, where $c$ and $d$ are
integers with $d\neq 0$.  Ask yourself this:  Are $a^y$ and $b^y$ necessarily integers?
\item
This is not correct.  If $y=3/2$, what does it mean to multiple $x$ by itself one and a half times?
\end{pfevalenum}
\end{answer}
\end{eval}

Since none of the proofs in the previous example were correct, you need to prove it.

\newpage
\begin{exer}
Prove or disprove that if $x$ and $y$ are rational, then $x^y$ is rational.\\
Proof:
\vspace*{2in}
\begin{answer}
The statement is false.  There are many counterexamples, but here is an easy one:
Let $x=2$ and $y=1/2$.  Then $x^y=2^{1/2}=\sqrt{2}$, which is irrational.
\end{answer}
\end{exer}

\begin{eval}
Prove or disprove that if $x$ is irrational, then $1/x$ is irrational.
\begin{spfenum}
\item
If $x$ is rational, assume it is an integer.  If $x$ is an integer, it is rational.
$1/x$ is an integer over an integer, so it is rational.  Therefore if $x$ is rational,
$1/x$ is rational, so by contrapositive reasoning, if $x$ is irrational, $1/x$ is irrational.\\
\evallinethree
\item
Assume that $x$ is irrational.  Then it cannot be expressed as an integer over an integer.
Then clearly $1/x$ cannot be expressed as an integer over an integer.\\
\evallinethree
\item
Assume that $x$ is rational.  Then $x=\frac{p}{q}$, where $p$ and $q$ are integers and $q\neq 0$.
But then $\displaystyle \frac{1}{x}=\frac{1}{\frac{p}{q}}=\frac{q}{p}$, so it is rational. 
Since we proved the contrapositive, the statement 
is true.\\
\evallinethree
\item
We will prove the contrapositive.  Assume that $1/x$ is rational.  
Since it is rational, $1/x=a/b$ for some integers
$a$ and $b$, with $b\neq 0$. Solving for $x$ we get $x=b/a$, so $x$ is rational.\\
\evallinethree
\item
I will prove the contrapositive statement:  If $1/x$ is
rational, then $x$ is rational.  Assume $1/x$ is rational.
Then $\frac{1}{x}=\frac{a}{b}$ for some integers $a$ and $b\neq 0$.
We know that $1/x\neq 0$ (since otherwise $x\cdot 0 = 1$, which is impossible), so
$a\neq 0$.  
Multiplying both sides of the previous equation by $x$ we get $x\frac{a}{b}=1$. Now if we multiply
both sides by $\frac{b}{a}$ (which we can do since $a\neq 0$), we get $x=\frac{b}{a}$.
Since $a$ and $b$ are integers with $a\neq 0$, $x$ is rational.\\
\evallinethree
\end{spfenum}
\begin{answer}
\begin{pfevalenum}
\item
This solution has two serious flaws.  First, we absolutely cannot assume $x$ is an integer.
The only thing we can assume about $x$ is that it is rational, and not every rational number
is an integer.  
The other problem is that the writer proved the {\em inverse}, not the {\em contrapositive}.
What they needed to prove was that if $1/x$ is rational, then $x$ is rational.  
So in actuality, we know is that $1/x$ is rational, not $x$.  We need to prove that $x$ is
rational based on the assumption that $1/x$ is rational.
\item
This is not really a proof.  It just takes the statement of the problem one step further.
Is the writer {\em sure} that $1/x$ can't be expressed as an integer over an integer?  Why?
There are just too many details omitted.
\item
The biggest flaw is that this is a proof of the {\em inverse} statement, not the {\em contrapositive}.
So even it the rest of the proof were correct, it would be proving the wrong thing
since the {\em inverse} and {\em contrapositive} are not equivalent.  But the rest is not
even entirely correct because the inverse statement is not quite true. 
If $x=0$, then $p=0$ as well and the statement and proof falls apart for the same reason--you can't divide by $0$.
\item
This proof is {\em almost correct}.  It does correctly try to prove the contrapositive, and if it had done so
correctly, that would imply the original statement is true.  
But there is one small problem:  If $a=0$ the proof
would fall apart because it would divide by $0$.  This possibility needs to be dealt with.  
This is actually not too difficult to fix.  We just need to add the following sentence before the last sentence:
{\em ``Since $0\neq 1/x$ for any value of $x$, we know that $a\neq 0$.''.}
\item
This proof is correct.
\end{pfevalenum}
\end{answer}
\end{eval}

%\newpage

\begin{eval}
Mersenne primes are primes that are of the form $2^p-1$, where $p$ is prime.
Are all numbers of this form prime?  Give a proof/counterexample.
\begin{spfenum}
\item
Restate the problem as if $2^p-1$ is prime then $p$ is prime.
Assume $p$ is not prime so $p=s t$, where $s$ and $t$ are integers.
Thus $2^p-1=2^{st}-1=(2^s-1)(2^{st-s}+2^{st-2s}+\cdots+2^s+1)$.
Because neither of these factors is $1$ or $2^p-1$\\
$\rightarrow$ $2^p-1$ is not prime (contradiction)\\
$\rightarrow$ $p$ is prime\\
$\rightarrow$ All numbers of the form $2^p-1$ (with $p$ a prime) are prime.\\
\evallinethree
\item
Numbers of the form $2^p$ only have $2$ as a factor.  Since $2^p-1$ is clearly odd, it does
not have $2$ as a factor.  Therefore it must not have any factors.  So it is prime.\\
\evallinetwo
\end{spfenum}
\begin{answer}
\begin{pfevalenum}
\item
As you will prove next, the statement is actually false.
Therefore the proof has to be incorrect.
But where did it go wrong?
It turns out they they tried to prove the wrong thing.  What needed to be proved was
``If $p$ is prime then $2^p-1$ is prime.''  They attempted to prove the {\em converse}
statement, which is not equivalent.  
We can still learn something by evaluating their proof.
It turns out that the converse is actually true, 
and the proof has a lot of correct elements.  Unfortunately, they are not put together
properly.   First of all, the proof seems to be a combination of a contradiction proof
and a proof by contrapositive.  They needed to pick one and stick with it.
Second, the arrows ($\rightarrow$) are confusing.  What do they mean?
I think they are supposed to be read as ``implies'', but a few more words are needed to
make the connections between these phrases.
Finally, the final statement is incorrect.  This does {\em not} prove that all
numbers of the form $2^p-1$ are prime when $p$ is prime.
\item
This proof is not even close.  
This is a case of ``I wasn't sure how to prove it so I just said stuff that sounded good.''
You can't argue anything about the factors of 
$2^p-1$ based on the factors of $2^p$.  Further, although $2^p-1$ being odd means
$2$ is not a factor, it doesn't tell us whether or not the number might have
{\em other} factors.
\end{pfevalenum}
\end{answer}
\end{eval}

\begin{exer}
Let $p$ be prime.  Prove that not all numbers of the form $2^p-1$ are prime.\\
Proof:
\vspace*{2in}
\begin{answer}
Notice that $11$ is prime but that $2^{11}-1=23\cdot 89$ is not.
Therefore, not all numbers of the form $2^p-1$, where $p$ is prime, are prime. 
\end{answer}
\end{exer}

